% !TEX program = xelatex
\documentclass[11pt,fontset=fandol]{ctexart}

\usepackage[a4paper,margin=1.9cm]{geometry}
\usepackage{booktabs}
\usepackage{tabularx}
\usepackage{array}
\usepackage{xcolor}
\usepackage{tikz}
\usepackage{hyperref}
\usetikzlibrary{arrows.meta,positioning,fit,backgrounds}

\definecolor{ink}{HTML}{17202A}
\definecolor{muted}{HTML}{5D6D7E}
\definecolor{lineblue}{HTML}{2F6FAB}
\definecolor{softblue}{HTML}{EAF3FF}
\definecolor{softgreen}{HTML}{EAF7EF}
\definecolor{softorange}{HTML}{FFF4E4}
\definecolor{softgray}{HTML}{F4F6F7}
\definecolor{accent}{HTML}{1F7A5B}

\hypersetup{
  colorlinks=true,
  linkcolor=lineblue,
  urlcolor=lineblue
}

\setlength{\parindent}{0pt}
\setlength{\parskip}{0.55em}
\renewcommand{\arraystretch}{1.2}
\newcommand{\code}[1]{\texttt{\detokenize{#1}}}
\newcolumntype{Y}{>{\raggedright\arraybackslash}X}
\newcolumntype{L}[1]{>{\raggedright\arraybackslash}p{#1}}

\tikzset{
  box/.style={
    rectangle,
    rounded corners=4pt,
    draw=#1!70,
    fill=#1!12,
    very thick,
    align=center,
    inner sep=6pt,
    minimum height=1.0cm
  },
  box/.default=lineblue,
  smallbox/.style={
    rectangle,
    rounded corners=3pt,
    draw=#1!70,
    fill=#1!10,
    thick,
    align=center,
    inner sep=5pt,
    minimum height=0.82cm
  },
  smallbox/.default=lineblue,
  arrow/.style={-{Latex[length=2.3mm]}, very thick, draw=muted},
  faint/.style={draw=muted!45, dashed, rounded corners=4pt}
}

\title{\bfseries OpenCollab 各模式的简单介绍}
\date{2026-06-29}

\begin{document}
\maketitle

\section{概述}

OpenCollab 的使用方式可以分成两层来看。第一层是任务怎么进入系统，包括交互式输入、普通 Eval、Workflow CLI，以及为 SWE-bench 写的专用 Eval 外壳。第二层是进入系统之后由谁决定下一步，包括单 Lead、YAML Team 和 Python Workflow。

普通 Eval 是一个很薄的批量任务外壳。它读取 JSONL，把每一行变成一个 \code{EvalTask}，给 OpenCollab 一个任务描述和一个可选运行环境，最后提取 \code{git diff} 写进结果文件。SWE-bench 专用 Eval 外壳更厚。它额外处理 SWE-bench 数据集、官方镜像、\code{/testbed} 容器、问题描述、隐藏测试字段、预测文件和官方评分器。两者都可以驱动 OpenCollab，但面向的任务协议不同。

\section{总图}

\begin{center}
\resizebox{0.98\linewidth}{!}{%
\begin{tikzpicture}[node distance=0.75cm and 1.55cm]
  \node[text=muted] (entrylabel) {入口层：任务以什么形态进入系统};
  \node[box=lineblue, minimum width=3.4cm, below=0.28cm of entrylabel] (interactive) {Interactive\\人工连续输入};
  \node[box=lineblue, minimum width=3.4cm, below=of interactive] (eval) {普通 Eval\\tasks.jsonl};
  \node[box=lineblue, minimum width=3.4cm, below=of eval] (workflowcli) {Workflow CLI\\--args JSON};
  \node[box=accent, minimum width=3.4cm, below=of workflowcli] (sweb) {SWE-bench 专用外壳\\dataset instance};

  \node[text=muted, right=3.2cm of entrylabel] (controllabel) {控制层：谁决定下一步};
  \node[smallbox=softgray, draw=muted, fill=softgray, minimum width=3.4cm, below=0.72cm of controllabel] (lead) {Single Lead\\LLM 自主调度};
  \node[smallbox=softgray, draw=muted, fill=softgray, minimum width=3.4cm, below=of lead] (team) {YAML Team\\角色与拓扑};
  \node[smallbox=softgray, draw=muted, fill=softgray, minimum width=3.4cm, below=of team] (pyflow) {Python Workflow\\代码定义流程};

  \node[box=softorange, draw=accent, fill=softorange, right=2.8cm of team, minimum width=4.6cm] (agent) {OpenCollab 执行层\\读文件、改代码、跑测试、生成 diff};
  \node[box=softgreen, draw=accent, fill=softgreen, below=0.85cm of agent, minimum width=3.4cm] (patch) {patch / prediction};
  \node[smallbox=softgreen, draw=accent, fill=softgreen, below=0.65cm of patch, minimum width=3.4cm] (judge) {官方或自定义评分};

  \draw[arrow] (interactive.east) -- (lead.west);
  \draw[arrow] (interactive.east) -- (team.west);

  \draw[arrow] (eval.east) -- (lead.west);
  \draw[arrow] (eval.east) -- (team.west);
  \draw[arrow] (eval.east) -- (pyflow.west);

  \draw[arrow] (workflowcli.east) -- (pyflow.west);

  \draw[arrow] (sweb.east) -- (lead.west);
  \draw[arrow] (sweb.east) -- (team.west);
  \draw[arrow] (sweb.east) -- (pyflow.west);

  \draw[arrow] (lead.east) -- (agent.west);
  \draw[arrow] (team.east) -- (agent.west);
  \draw[arrow] (pyflow.east) -- (agent.west);
  \draw[arrow] (agent) -- (patch);
  \draw[arrow] (patch) -- (judge);

  \begin{scope}[on background layer]
    \node[faint, fit=(interactive)(eval)(workflowcli)(sweb)(entrylabel)] {};
    \node[faint, fit=(lead)(team)(pyflow)(controllabel)] {};
    \node[faint, fit=(agent)(patch)(judge)] {};
  \end{scope}
\end{tikzpicture}
}
\end{center}

\section{普通 Eval 外壳}

普通 Eval 的命令形态是 \code{opencollab eval tasks.jsonl}。这个 JSONL 文件每行代表一个任务，核心字段是 \code{task_id} 和 \code{description}，可选字段包括 \code{repo_path}、\code{docker_image}、\code{timeout} 和 \code{max_tokens}。CLI 读入这些字段之后构造 \code{EvalTask}，再交给通用 evaluator 执行。

\begin{center}
\resizebox{0.94\linewidth}{!}{%
\begin{tikzpicture}[node distance=1cm]
  \node[box=lineblue, minimum width=3.1cm] (jsonl) {tasks.jsonl\\任务描述};
  \node[box=lineblue, right=of jsonl, minimum width=3.1cm] (task) {EvalTask\\统一任务对象};
  \node[box=lineblue, right=of task, minimum width=3.4cm] (env) {Local 或 Docker\\运行环境};
  \node[box=lineblue, right=of env, minimum width=3.3cm] (agent) {OpenCollab\\单次会话};
  \node[box=lineblue, right=of agent, minimum width=3.0cm] (result) {results.jsonl\\patch 与轨迹};
  \draw[arrow] (jsonl) -- (task);
  \draw[arrow] (task) -- (env);
  \draw[arrow] (env) -- (agent);
  \draw[arrow] (agent) -- (result);
\end{tikzpicture}
}
\end{center}

普通 Eval 的职责边界很清楚。它负责批量读取任务、准备本地或 Docker 环境、发送一次初始任务描述、运行 OpenCollab、提取最终 diff、保存通用结果。它不理解 SWE-bench 的官方镜像命名，不解析 \code{problem_statement}，不生成 SWE-bench 官方 \code{predictions.jsonl}，也不直接调用官方评分器。

\section{SWE-bench 专用 Eval 外壳}

SWE-bench 专用外壳的输入不是普通 JSONL 任务，而是 SWE-bench 数据集实例。每个实例包含 \code{instance_id}、\code{repo}、\code{problem_statement}、\code{hints_text}、\code{FAIL_TO_PASS} 等 benchmark 字段。外壳先用 SWE-bench 官方工具确定实例镜像，再在官方 \code{/testbed} 环境中启动 OpenCollab。

\begin{center}
\resizebox{0.98\linewidth}{!}{%
\begin{tikzpicture}[node distance=0.85cm and 0.95cm]
  \node[box=accent, minimum width=3.2cm] (dataset) {SWE-bench\\dataset instance};
  \node[box=accent, right=of dataset, minimum width=3.0cm] (spec) {make\_test\_spec\\官方镜像名};
  \node[box=accent, right=of spec, minimum width=3.4cm] (container) {sweb.eval.*\\ /testbed 容器};

  \node[box=lineblue, below=1.1cm of dataset, minimum width=3.2cm] (prompt) {problem statement\\hints / F2P};
  \node[box=lineblue, right=of prompt, minimum width=3.0cm] (oc) {OpenCollab\\Lead / Team / Workflow};
  \node[box=lineblue, right=of oc, minimum width=3.4cm] (diff) {git diff\\model\_patch};

  \node[box=softgreen, draw=accent, fill=softgreen, below=1.1cm of oc, minimum width=3.2cm] (pred) {predictions.jsonl};
  \node[box=softgreen, draw=accent, fill=softgreen, right=of pred, minimum width=3.4cm] (official) {official harness\\resolved / failed};

  \draw[arrow] (dataset) -- (spec);
  \draw[arrow] (spec) -- (container);
  \draw[arrow] (dataset) -- (prompt);
  \draw[arrow] (container) -- (oc);
  \draw[arrow] (prompt) -- (oc);
  \draw[arrow] (oc) -- (diff);
  \draw[arrow] (diff) -- (pred);
  \draw[arrow] (pred) -- (official);
\end{tikzpicture}
}
\end{center}

这套外壳做的是 benchmark 适配。它把 SWE-bench 自己的任务格式、镜像格式、评测格式转换成 OpenCollab 的一次解题过程，再把 OpenCollab 产生的 diff 转回 SWE-bench 官方预测格式。我们之前主要使用的 Team 评测路径，就是由 \code{scripts/run_team_batch.sh} 选择实例并解析镜像，再由 \code{scripts/start_team_run.sh} 启动容器、构造 prompt、调用 OpenCollab。

\end{document}
